% !TeX root = ../../../金融经济学_王江_讲义及习题.tex
% 来源：教材图12.2。用可行组合 q、q' 展示候选射线斜率逐步增大。
% 纯风险前沿 sigma=sqrt(1+t^2), r=1.8+.7t，r_F=.9。
% 切点 t=7/9，切线斜率 sqrt(1.3)。
\begin{tikzpicture}[>=Stealth,x=1.6cm,y=1.05cm,font=\small]
 \draw[->] (0,0)--(3.7,0) node[right] {$\sigma$};
 \draw[->] (0,0)--(0,4.75) node[above] {$\bar r$};
 \node[below left] at (0,0) {$O$};
 \draw[black!80,thick,domain=-2.8:3.15,samples=110,variable=\t]
   plot ({sqrt(1+\t*\t)},{1.8+.7*\t});
 \draw[gray!65,densely dashed] (0,.9)--(3.3,.9);
 \draw[gray!85,dashed] (0,.9)--(3.25,2.40893);
 \draw[gray!85,dashed] (0,.9)--(3.25,3.10536);
 \draw[blue!70!black,very thick] (0,.9)--(3.2,4.54856);
 \node[blue!70!black,above left] at (3.2,4.54856) {最大 Sharpe 比};
 \draw[->,gray!85] (.9,.9) arc[start angle=0,end angle=48.75,radius=.9];
 \node[anchor=west,font=\footnotesize] at (1.05,1.15) {Sharpe 比增大};
 \fill (0,.9) circle (1.8pt) node[left] {$r_F$};
 \filldraw[fill=white] (2.8,2.2) circle (2pt) node[below right=2pt] {$q$};
 \filldraw[fill=white] (2.8,2.8) circle (2pt) node[below right=2pt] {$q'$};
 \filldraw[draw=blue!70!black,fill=white] (1.26686,2.34444)
   circle (2.2pt) node[above left=3pt,black] {$T$};
\end{tikzpicture}
